% 若编译失败，且生成 .synctex(busy) 辅助文件，可能有两个原因：
% 1. 需要插入的图片不存在：Ctrl + F 搜索 'figure' 将这些代码注释/删除掉即可
% 2. 路径/文件名含中文或空格：更改路径/文件名即可

% --------------------- 文章宏包及相关设置 --------------------- %
% >> ------------------ 文章宏包及相关设置 ------------------ << %
% 设定文章类型与编码格式
\documentclass[UTF8]{article}		
\input{../../.config/config_for_LinearCircuitExperiment.tex}



%%%%%%%%%%%%%%%%%%%%%%%%%%%%%%%%%%%%%%%%%%%%%%%%%%%%%%%%%%%%%%%%
% 仅需修改页眉、实验名称、实验日期
%%%%%%%%%%%%%%%%%%%%%%%%%%%%%%%%%%%%%%%%%%%%%%%%%%%%%%%%%%%%%%%%


%%%%%%%%%%%%%%%%%% 1. 修改页眉内容 %%%%%%%%%%%%%%%%%%
\rhead{Preview Report of LCE-08 运放参数测量 (2025.05.30, 丁毅)}

% 开始编辑文章
\begin{document}
\begin{center}\large
    \vspace*{-0.8cm}
    \noindent{\huge\bfseries《\ \ 线\ \ 性\ \ 电\ \ 路\ \ 实\ \ 验\ \ \ 》\ \ 预\ \ 习\ \ 报\ \ 告 }
    \\\vspace{0.1cm}
    \noindent{
    {\bfseries 
%
%%%%%%%%%%%%%%%%%% 2. 修改实验名称 %%%%%%%%%%%%%%%%%%
    实验名称：\uline{\hspace{1.5cm} 运放参数测量 \hspace{1.5cm}}
%
    }\hspace{0.4cm}
    指导教师：\uline{\hspace{0.8cm}王东雷\ \ \ \  \ df4dac@sina.com     \hspace{0.8cm}}
    }
    \\\vspace{0.1cm}
    \noindent
    {
    姓名：\uline{\,\,\,丁毅\,\,\,}\hspace{0.2cm}
    学号：\uline{\,\,\,{ 2023K8009908031}\,\,\,}\hspace{0.2cm}
    班级/专业：\uline{\,\,\,{2308/电子信息}\,\,\,}\hspace{0.2cm}
    分组序号：\uline{\,\,\,{2-06}\,\,\,}
    }
    \\\vspace{0.1cm}
    \noindent{
%
%%%%%%%%%%%%%%%%%% 3. 修改实验日期 %%%%%%%%%%%%%%%%%%
    实验日期：\uline{\,\,{2025.05.30}\,\,}\hspace{0.2cm}
%
    实验地点：\uline{\,\,\,教学楼{ 607}\,\,\,}\hspace{0.2cm}
    是否调课/补课：\uline{\hspace{0.26cm}否 \hspace{0.26cm}}\hspace{0.2cm}
    成绩：\uline{\hspace{2cm}}
    }
\end{center}
\vspace{-0.4cm}
\noindent\rule{\textwidth}{0.075em}   % 分割线
\vspace{-1.0cm}


% ------------------------ 文章信息区 ------------------------ %
% ------------------------ 文章信息区 ------------------------ %



%%%%%%%%%%%%%%%%%%%%%%%%%%%%%%%%%%%%%%%%%%%%%%%%%%%%%%%%%%%%%%%%%%%%%%%%%%%%%%%%%
%%%%%%%%%%%%%%%%%%%%%%%%%%%%%%%%% 下面是正文内容 %%%%%%%%%%%%%%%%%%%%%%%%%%%%%%%%%
%%%%%%%%%%%%%%%%%%%%%%%%%%%%%%%%%%%%%%%%%%%%%%%%%%%%%%%%%%%%%%%%%%%%%%%%%%%%%%%%%

\section{实验目的}

\begin{enumerate}
\item 理解运放主要参数的定义及对电路性能的影响，增强对运放应用电路的理解；
\item 掌握运放主要参数的测量方法；
\item 测量 uA741, LM358, LM318N, OP177 等典型集成运放的基本参数
\end{enumerate}



\section{实验仪器}

\begin{enumerate}
    \item 数字万用表： Unit UT61E (C190241394)
    \item 数字示波器： RIGOL 200MSO2202A (DS2F192200361)
    \item 信号发生器： GWINSTEK AFG-22225 (GER910370)
    \item 数字直流电源： GWINSTEK GPD-3303S (GES813705) 
    \item 多功能数字测量仪： 
    % 这里是链接
    \href{https://digilent.com/reference/test-and-measurement/analog-discovery/start
    }{ % 这里是文字
    Analog Discovery 1
    } 
    (D704387)
    \item 运放基本参数测试板：% 这里是链接  
    \href{https://yidingg.github.io/YiDingg/\#/ElectronicDesigns/Basic\%20Op\%20Amp\%20Measurement\%20Board\%20v2
    }{ % 这里是文字
    Basic Op Amp Measurement Board v2
    }
    \item 电压线性运算器：
    % 这里是链接
    \href{https://yidingg.github.io/YiDingg/#/ElectronicDesigns/Voltage\%20Linear\%20Operation\%20Board
    }{ % 这里是文字
    Voltage Linear Operaion Board
    }
    \item 其它：运算放大器 uA741, LM258P, LM318N, OP177 和 OPA197、面包板，电容、电阻、杜邦线等
\end{enumerate}

\section{实验内容概要}


注：我们并不按照讲义上给出的具体步骤进行实验，所有数据均以图片形式先记录在个人网站中，再统一转到 MATLAB 进行处理，因此不需要纸质数据表格。
\begin{enumerate}
\item 测量 uA741 的全部基本参数，包括 $V_{OS},\ I_{B\pm}$, DC 和 AC 下的 $A_v,\ \mathrm{PSRR},\ \mathrm{CMRR}$, slew rate 和 FPBW
\item 测量 LM258P 的 $V_{OS},\ I_{B\pm}$
\item 测量 OP177 的 $V_{OS},\ I_{B\pm}$ 和 $A_{v,DC}$
\item 测量 LM318N 的 slew rate, 以及在 10 Vamp output 下的输出摆幅曲线 (理论带宽 $\displaystyle f_p = \frac{\mathrm{SR}}{2\pi V_{out,amp}}$)
\item (选做) 测量 OPA197 的各项参数，包括 $V_{OS},\ I_{B\pm}$, DC 和 AC 下的 $A_v,\ \mathrm{PSRR},\ \mathrm{CMRR}$ 以及 slew rate
\end{enumerate}

\section{uA741 原理图}

\begin{figure}[H]\centering
    \includegraphics[width=\columnwidth]{assets/uA741_schematic.pdf}
    \caption{Circuit schematic of uA741}
\end{figure}




%\section{异常现象分析}
















































\end{document}

% VScode 常用快捷键：

% F2:                       变量重命名
% Ctrl + Enter:             行中换行
% Alt + up/down:            上下移行
% 鼠标中键 + 移动:           快速多光标
% Shift + Alt + up/down:    上下复制
% Ctrl + left/right:        左右跳单词
% Ctrl + Backspace/Delete:  左右删单词    
% Shift + Delete:           删除此行
% Ctrl + J:                 打开 VScode 下栏(输出栏)
% Ctrl + B:                 打开 VScode 左栏(目录栏)
% Ctrl + `:                 打开 VScode 终端栏
% Ctrl + 0:                 定位文件
% Ctrl + Tab:               切换已打开的文件(切标签)
% Ctrl + Shift + P:         打开全局命令(设置)

% Latex 常用快捷键：

% Ctrl + Alt + J:           由代码定位到PDF


